\documentclass[../Main.tex]{subfiles}
\begin{document}

\section{Chosen Technologies}
\label{chap:chosen_tech}
Deciding which technologies we will use when developing our chosen game concept, Hover, is an essential consideration in the project. In this chapter, a description of the selection process and the chosen technologies will be presented. 

\subsection{Selection Process}
Chapter \ref{chap:technology_review} presents an overview of relevant technology. The theory presented here will be used as a basis for selecting technology for developing our concept, Hover. Because Hover is based on the mobile platform, there are essentially two main technology choices we need to undertake: Decide which mobile application development approach and which backend solution approach we want to use. After some discussion, we have decided to choose a cross-platform framework for developing the mobile application and a Backend as a Service (BaaS) solution for managing data, users, and analytics.

Our main reason for choosing a cross-platform framework instead of the native approach for mobile development is justified based on the possible target group. As explained earlier, in our motivation for doing this project, we want to create a game that can reach as many users as possible. By choosing a cross-platform mobile framework, we can create an application that reaches users on both Android and iOS. Additionally, developing a game that runs on Android and iOS increases the number of possible test participants. Making it easier for us to construct an appropriate test group when we are going to evaluate our concept. As long as the framework exposes the possibilities we need, there are little downsides to using such a tool.

When choosing our backend solution approach, the main concern was simplicity. We want to maximize the time spent on developing the actual game and increase the functionality available. It is important for us that the realization of the concept is as complete as possible before testing the game, to get a good evaluation of it. The development time in the master thesis period is limited. Thus, spending a lot of time on a complex backend architecture may be less beneficial for creating a good evaluation of the game concept.

After choosing the main mobile development and backend approach, we turn our focus to choosing specific tools we want to utilize. When selecting the cross-platform mobile application framework and the BaaS solution, we used three main evaluation criteria:

\textbf{Maturity:} It is essential for us that the tool we choose is mature. A tool that has been tested,  experienced, and used in many applications is positive. It is also essential that the solution is undergoing active development and maintenance. Open source is an additional plus, but not a requirement. Open source tools often lead to more active development since in principle everyone can contribute to the project.

\textbf{Support for required functionality:}
An essential aspect is that central functionality is possible to develop. The chosen tools need to have support for what we want to do. Specifically, in our concept, this includes location tracking, map integration as well as authentication, and database for the backend solution.

\textbf{Previous experience:} Prior knowledge of a tool counts positive because it will increase the development speed. While it is always fun to try out new technologies as a developer, some previous experience is preferable. The main goal of this project is to invent, develop, and test a game concept, and thus it is important that the game is well implemented.
\subsection{Description of Technologies}
Based on the selection criteria described above we have selected React Native as a mobile application framework and Firebase as a BaaS provider. 
\subsubsection{React Native as a Mobile Development Framework}
React Native is a cross-platform framework developed by Facebook introduced at React.js Conf, in 2015 \citep{occhino_react_2015}. The main goal of the framework is to be able to write applications that can run on both iOS and Android. One of the main differences between React Native and other cross-platform frameworks is its use of native components. For example, a button declared in a React Native application will be presented as an Android-type button when the application is run on Android, and an iOS-type button when it is used on iOS. Most other frameworks have a different approach where the application provides a common experience for both the Android and the iOS application. In this way, a React Native application feels more like a native application than other frameworks such as Flutter and Ionic. As described in the section above we had three main criteria when choosing React Native, namely, maturity, support for required functionality, and previous experience.

\textbf{Maturity:} React Native is a mature mobile development framework and it is used by a handful of Fortune-500 \href{https://reactnative.dev/showcase}{companies} for their mobile applications. The framework is open-source and under active development and maintenance. Because React Native is based on React and JavaScript, the framework also has access to NPM’s large collection of packages of open-source code.

\textbf{Support for required functionality:} Location tracking, geofencing, and map-view areas mentioned earlier are essential functionality in Hover. Fortunately, there exist libraries for all these functionalities in React Native.

\textbf{Previous experience:} We do not have a lot of experience with mobile application development, except for some limited experimentation and prototyping in conjunction with the technology review process. However, we have some experience working with React and this is a definite plus when developing in React Native.

Although we chose React Native as our development framework, other frameworks such as Flutter and Xamarin would presumably also have been great options. In summary, it was a combination of small differences, where our previous experience with React was the deciding factor that determined our selection.

\subsubsection{Firebase as Backend as a Service}
Firebase is a Backend as a Service solution provided by Google. The service was launched as a separate company back in 2011 and bought by Google in 2014 \citep{tamplin_firebase_2014}. Firebase includes tools for storage, analytics, cloud functions, authentication, and more. Its main selling point is the simplicity of using it accompanying the development of a mobile application. Similar to when choosing the front-end tool, the three main evaluation criteria were also for our decision to use Firebase:

\textbf{Maturity:} Released back in 2011, Firebase was one of the first backend solutions designed specifically to support mobile applications \citep{metz_firebase_2012}. The solution is used by \href{https://firebase.google.com/use-cases}{companies} across the world and is consequently highly maintained and developed. Firebase also develops open-source libraries and extensions for use in various frameworks and languages. 

\textbf{Support for required functionality:} Firebase has support for the backend tools needed to develop Hover. This is primarily authentication, storage, and analytics. Authentication will be used to verify users, storage for saving game data, while analytics will be used to see how the users are utilizing the application during the test period.

\textbf{Previous experience:} Firebase was the only BaaS we had some prior knowledge of. Even though our experience is not extensive, we do have experience with some of the fundamental parts of Firebase. 

The criteria valued the most in this decision was the previous experience. Both of us had some prior knowledge of Firebase. After doing some research it seems like there are also more tools created for Firebase in React Native than any of the other services providers. Even though we chose Firebase, AWS Amplify or Azure Mobile apps could also have been good alternatives.

\subsubsection{Testing the technologies in a quick prototype}
To make sure that the connection between Firebase and React Native, and other related important libraries we are going to use, are working as expected, we have developed a prototype in parallel with writing this chapter. The prototype includes a mobile application written in React Native, communicating with a database in Firebase. We have implemented a simple map-view centered on the user’s location and fetching of dummy data from storage in Firebase. More complex functionality such as authentication, geo-fencing, and game ideas has not been prioritized. In this way, we have tested the most essential parts required for the game setup. Figure ?? and Figure ?? show screenshots of the map-view and the dummy data view from the prototype.

\textbf{[Insert screenshots]}

\subsection{Summary}
We have looked at the chosen technologies and some considerations regarding these technologies. After discussions and a lot of research and experimentation, we ended up choosing a combination of React Native and Firebase. In addition to selecting the technologies, we have developed a quick prototype to ensure the most essential functionality is working as expected. The choices and research we have done here prepare us for the development period in the master thesis period.

\biblio % Needed for referencing to working when compiling individual subfiles - Do not remove
\end{document}